\section{Proofs for Metric and Split-Hadamard Reduction}
\label{app:metric-proofs}

\begin{proof}[Proof of \cref{thm:metric-characterization}]
Assume (A).  If \(z\in X\) and \(r=d_X(x,z)\), ball surjectivity implies
\(d_Y(q(x),q(z))\le d_X(x,z)\); hence \(q\) is \(1\)-Lipschitz.  For fixed
\(y\), set \(\delta=d_Y(q(x),y)\).  Since
\(y\in\overline B_Y(q(x),\delta)\), submetry gives
\(z\in\overline B_X(x,\delta)\) with \(q(z)=y\).  Contraction yields
\[
  \delta\le d_X(x,z)\le\delta,
\]
proving (B), including attainment.

Assume (B).  Taking \(y=q(z)\) gives
\[
  d_Y(q(x),q(z))
  =\min_{q(w)=q(z)}d_X(x,w)
  \le d_X(x,z),
\]
so \(q\) is \(1\)-Lipschitz and maps the ambient ball into the visible ball.
If \(y\) lies in the visible ball, an attaining fiber minimizer lies in the
ambient ball.  Thus (B) implies (A).

Assume (B), fix \(z\), and put \(y=q(z)\).  Then
\(d_X(x,z)\ge d_Y(q(x),y)\), so monotonicity of \(\rho\) gives the lower
bound from the left side of \eqref{eq:universal-radial-reduction} to the
right.  An attaining fiber minimizer realizes equality for every \(y\),
proving the reverse inequality and hence (C).

Finally assume (C).  For \(r\ge0\) and fixed \(y\), use
\[
  \rho_r(t)=\begin{cases}0,&t\le r,\\+\infty,&t>r,\end{cases}
  \qquad
  \phi_y(u)=\begin{cases}0,&u=y,\\+\infty,&u\ne y.\end{cases}
\]
The reduced infimum is zero exactly when
\(y\in\overline B_Y(q(x),r)\), whereas the ambient infimum is zero exactly
when a point of \(q^{-1}(y)\) lies in \(\overline B_X(x,r)\).  Equality for
all \(y,r\) is precisely the submetry identity.  Thus (C) implies (A).
\end{proof}

\begin{proof}[Proof of \cref{thm:coherent-reduction}]
Assume (A).  If \(q(z)=y\), contraction gives
\[
  d_X(x,z)\ge d_Y(q(x),y).
\]
Because \(x=s_x(q(x))\) and \(s_x\) is isometric,
\[
  d_X(x,s_x(y))=d_Y(q(x),y),
\]
which proves (B).

Assume (B), fix \(y_1,y_2\), and put \(x_1=s_x(y_1)\).  Coherence gives
\(s_{x_1}=s_x\).  Applying (B) with anchor \(x_1\) and target \(y_2\),
\[
  d_X(s_x(y_1),s_x(y_2))
  =d_Y(q(x_1),y_2)=d_Y(y_1,y_2),
\]
so (B) implies (A).

Under (B), every \(z\) with \(q(z)=y\) has
\[
  \rho(d_X(x,z))+\phi(y)
  \ge \rho(d_Y(q(x),y))+\phi(y),
\]
and \(z=s_x(y)\) realizes equality.  Taking infima proves universal
reduction and canonical minimizer lifting, so (B) implies (C).

Assume (C), choose \(\rho(t)=t\), and let \(\phi\) be the indicator of a
singleton \(\{y\}\).  Universal reduction identifies the fiber infimum with
\(d_Y(q(x),y)\), while the canonical-minimizer clause says that \(s_x(y)\)
attains it.  Hence (C) implies (B).

The isometric-section property gives the reverse ball inclusion exactly as in
the proof of \cref{thm:metric-characterization}; hence \(q\) is a submetry.
For the proximal identity, apply universal reduction with
\(\rho(t)=t^2/(2\tau)\).  A proper lower-semicontinuous geodesically convex
function plus squared distance has a unique minimizer in complete CAT(0)
space.  The section lifts it with equal value.  Coherence then gives exact
iteration by induction.
\end{proof}

\section{Proof of the Split-Hadamard Theorem}
\label{app:split-proofs}

\begin{lemma}[Complete local isometries over a Hadamard base]
\label{lem:complete-local-isometry}
Let \(F:L\to B\) be a local Riemannian isometry from a connected complete
Riemannian manifold to a Hadamard manifold.  Then \(F\) is a global
Riemannian isometry.
\end{lemma}

\begin{proof}
Fix \(p\in L\).  Any geodesic \(\gamma:[0,1]\to B\) beginning at \(F(p)\)
has a unique local lift beginning at \(p\), obtained from a local inverse of
\(F\).  A lifted segment has the same speed as \(\gamma\).  If its maximal
interval ended at \(T\le1\), then for \(s,t\uparrow T\),
\[
  d_L(\widetilde\gamma_s,\widetilde\gamma_t)
  \le \operatorname{Len}(\widetilde\gamma|_{[s,t]})
  =\operatorname{Len}(\gamma|_{[s,t]})\longrightarrow0.
\]
Completeness gives a limit point in \(L\), and a local inverse at that point
extends the lift past \(T\).  Thus every such geodesic lifts over its entire
interval.  Since every point of the connected complete manifold \(B\) is
joined to \(F(p)\) by a minimizing geodesic, \(F\) is surjective.

The same continuation argument gives unique lifting for every piecewise
smooth path and for homotopies through such paths; equivalently, a complete
local isometry is a covering map.  The Hadamard manifold \(B\) is simply
connected, so its connected covering \(L\) has one sheet.  Hence \(F\) is a
diffeomorphism.  A bijective local Riemannian isometry preserves lengths in
both directions and is therefore a global Riemannian isometry.
\end{proof}

\begin{proof}[Proof of \cref{thm:split-hadamard}]
Because \(L_x\) is a horizontal integral leaf,
\(Dq|_{TL_x}\) is an isometry at every point, so \(q|_{L_x}\) is a local
Riemannian isometry.  The leaf is connected and complete by definition, and
the Hadamard base is simply connected.  Applying
\cref{lem:complete-local-isometry} shows directly that
\(q|_{L_x}:L_x\to B\) is a global Riemannian isometry.

The complete totally geodesic leaf \(L_x\) is geodesically convex in \(M\).
Indeed, Hopf--Rinow gives an intrinsic minimizing geodesic between any two
points of the leaf.  Total geodesy makes it an ambient geodesic, and the
Hadamard property makes that ambient geodesic the unique globally minimizing
one.  In particular, intrinsic and ambient distances agree on \(L_x\).
Any horizontal section through \(x\) lies in the same maximal integral leaf
and equals the inverse \(s_x\) of this global isometry.

For every piecewise smooth ambient curve \(\eta\), Riemannian-submersion
contraction gives
\[
  \operatorname{Len}(q\circ\eta)\le\operatorname{Len}(\eta),
\]
and hence \(d_B(q(y),q(z))\le d_M(y,z)\).  If \(q(z)=b\),
\[
  d_M(x,z)\ge d_B(q(x),b),
\]
while \(s_x(b)\) attains equality.  Suppose another \(z\) also attains it,
and let \(\eta\) be the ambient minimizing geodesic from \(x\) to \(z\).
Every inequality in
\[
  d_B(q(x),b)
  \le\operatorname{Len}(q\circ\eta)
  \le\operatorname{Len}(\eta)
  =d_M(x,z)
\]
is an equality.  The orthogonal decomposition
\(\dot\eta=\dot\eta^{\mathcal H}+\dot\eta^{\mathcal V}\) then forces
\(\dot\eta^{\mathcal V}=0\) everywhere.  The geodesic stays in \(L_x\) and
ends at \(s_x(b)\), proving uniqueness.  The ball identity and radial
reduction now follow from \cref{thm:metric-characterization}.

Proximal commutation follows from radial reduction with squared distance.  The
reduced objective has a unique minimizer on the Hadamard base.  Equality in
the reduction first forces the compression of any full minimizer to equal that
visible proximal point and then forces the full point to be its unique shortest
lift from \(x\).  Thus the full proximal point is unique even though no
geodesic-convexity assertion for \(\phi\circ q\) on \(M\) is needed.  Finally, for
\(v\in T_zM\),
\[
  \dd(\phi\circ q)_z[v]
  =h_{q(z)}(\operatorname{grad}_B\phi(q(z)),Dq(z)[v]).
\]
This vanishes on \(\mathcal V_z\); on \(\mathcal H_z\), \(Dq\) is an
isometry.  Riesz representation gives \eqref{eq:split-gradient}.  Since
\(Ds_x\) is the inverse horizontal isometry, the chain rule shows that
\(z=s_x\circ y\) solves the lifted equation on every interval on which the
visible solution \(y\) exists.  Local Lipschitz continuity gives uniqueness
of maximal solutions by the standard ODE theorem.
\end{proof}

\begin{proof}[Proof of \cref{cor:split-coherent}]
The sections are isometric.  If \(z=s_x(y)\), then \(z\) and \(x\) lie on
the same maximal horizontal leaf, so both \(s_z\) and \(s_x\) invert the
restriction of \(q\) to that leaf.  Thus \(s_z=s_x\), proving coherence.
The unique shortest-lift statement gives rigidity.
\end{proof}
